\chapter{द्विविमीय शुद्धगतिकी}\label{ch:g12:kinematics-2d}

किसी कार को गोल चक्कर पर स्थिर \qty{30}{km/h} से घूमते देखिए: चालमापी की सुई हिलती तक
नहीं, फिर भी हर सवारी को बग़ल की ओर खिंचाव महसूस होता है --- यानी गति बदल रही
है, और बदल रही चीज़ चाल नहीं है। अब हर क्षण के लिए एक संख्या काफ़ी नहीं
रही: समतल में गति के लिए तीर चाहिए। यह अध्याय चलते बिंदु को तीन तीर सौंपता है
--- स्थिति, वेग, त्वरण --- जिन्हें अवकलज एक शृंखला में बाँधता है, और फिर गति
के पूरे-के-पूरे कुल उनकी ज्यामिति से पढ़ लेता है।

\section{स्थिति सदिश}

गति केवल किसी स्पष्ट रूप से बताए गए \omterm{def:g10:relative-motion:frame}{निर्देश तंत्र} के सापेक्ष होती है ---
यानी एक दृढ़ पिंड और एक घड़ी (\cref{ch:g10:relative-motion})। इस वर्ष नई बात परिशुद्धता है: उस तंत्र
में एक मूल बिंदु $O$ और दो परस्पर लंब अंशांकित अक्ष $Ox$, $Oy$ जोड़ दीजिए,
और हर क्षण बिंदु कहाँ है यह फलनों के एक युग्म के रूप में दर्ज कीजिए।

\begin{definition}[स्थिति सदिश]\label{def:g12:kinematics-2d:position}
मूल बिंदु $O$ और अक्षों $Ox$, $Oy$ से सजे तंत्र में समय $t$ पर चलते
बिंदु $M$ का \emph{स्थिति सदिश}\index{स्थिति सदिश} $\vect{OM}(t)$ है, जिसके
निर्देशांक $\bigl(x(t),\, y(t)\bigr)$ हैं: यानी समय के दो फलन, हर अक्ष के लिए एक। $t$ के बदलते
जाने पर $M$ जो वक्र खींचता है वही उसका \omterm{def:g10:relative-motion:trajectory}{गतिपथ} है।
\end{definition}

\section{वेग सदिश}

\begin{definition}[वेग सदिश]\label{def:g12:kinematics-2d:velocity}
किसी छोटे अंतराल में बिंदु $\vect{OM}(t+\Delta t) - \vect{OM}(t)$ जितना खिसकता है; उसे $\Delta t$ से भाग दीजिए और
अंतराल को सिकुड़ने दीजिए: यही आपके गणित के पाठ्यक्रम का अवकलज है, हर निर्देशांक
पर लगाया हुआ। $M$ का \emph{वेग सदिश}\index{वेग सदिश} उसके \omterm{def:g12:kinematics-2d:position}{स्थिति
सदिश} का अवकलज है, $\vect v(t) = \bigl(x'(t),\, y'(t)\bigr)$, जिसका मात्रक \unit{m/s} है। उसका मानक $v = \sqrt{x'^2 + y'^2}$ ही
\emph{चाल}\index{चाल} है --- वही एक संख्या जो चालमापी दिखाता है।
\end{definition}

\begin{example}[समतल उड़ान भरता ड्रोन]\label{ex:g12:kinematics-2d:drone}
कोई ड्रोन $x(t) = 4.0\,t$ और $y(t) = 3.0\,t$ (\omterm{def:g10:orders-of-magnitude:unit}{मीटर}, \omterm{def:g10:orders-of-magnitude:unit}{सेकंड}) के साथ उड़ता है: \omterm{def:g10:relative-motion:trajectory}{गतिपथ} रेखा
$y = \tfrac34 x$, वेग हर समय $\vect v = (4.0, 3.0)$, और चाल $\sqrt{4.0^2 + 3.0^2} = \qty{5.0}{m/s}$, अचर।
\end{example}

\begin{proposition}[गतिपथ का स्पर्शी]\label{prop:g12:kinematics-2d:tangent}
हर क्षण $\vect v(t)$ $M$ की वर्तमान स्थिति पर \omterm{def:g10:relative-motion:trajectory}{गतिपथ} का स्पर्शी होता है,
और चलने की दिशा में संकेत करता है।
\end{proposition}

\begin{proof}
विस्थापन \omterm{def:g10:relative-motion:trajectory}{गतिपथ} की एक जीवा है; और $\Delta t$ के सिकुड़ते ही जीवा की दिशा
स्पर्शी की दिशा बन जाती है।
\end{proof}

\section{त्वरण सदिश}

\begin{definition}[त्वरण सदिश]\label{def:g12:kinematics-2d:acceleration}
$M$ का \emph{त्वरण सदिश}\index{त्वरण सदिश} उसके \omterm{def:g12:kinematics-2d:velocity}{वेग सदिश} का अवकलज
है, $\vect a(t) = \bigl(v_x'(t),\, v_y'(t)\bigr) =
\bigl(x''(t),\, y''(t)\bigr)$, जिसका मात्रक \unit{m/s^2} है: यानी \omterm{def:g12:kinematics-2d:velocity}{वेग सदिश} इस समय कितनी तेज़ी
से और किस दिशा में बदल रहा है।
\end{definition}

\begin{remark}[चाल बढ़ाए बिना त्वरण]\label{rem:g12:kinematics-2d:turning}
जब भी वेग-\emph{सदिश} बदलता है --- लंबाई में \emph{या दिशा में} --- $\vect a \neq \vect 0$:
ब्रेक लगाती कार त्वरित होती है (पीछे की ओर); अचर चाल पर मुड़ती कार
त्वरित होती है (बग़ल की ओर)। यही वह राशि है जिसे बल नियंत्रित करते हैं
(\cref{ch:g11:forces-and-motion,ch:g12:newtons-laws})।
\end{remark}

\begin{method}[कालक्रम-चित्र पढ़ना]\label{met:g12:kinematics-2d:chrono}
\emph{कालक्रम-चित्र}\index{कालक्रम-चित्र} किसी चलते बिंदु की स्थितियाँ
$M_0, M_1, M_2, \dots$ बराबर समय-अंतरालों $\Delta t$ पर दर्ज करता है।
\begin{enumerate}
  \item बिंदुओं का अंतराल चाल बताता है: बराबर अंतराल, अचर
        चाल; खुलते अंतराल, चाल बढ़ रही है; सिकुड़ते अंतराल, घट रही
        है।
  \item $M_i$ पर वेग सममित जीवा है, $\vect v_i \approx \vect{M_{i-1}M_{i+1}}/(2\,\Delta t)$, जो $M_i$ पर खींची जाती है और
        बिंदुओं वाले पथ की स्पर्शी होती है।
  \item $M_i$ पर त्वरण के लिए: $\vect v_{i-1}$ और $\vect v_{i+1}$ को $M_i$ पर उतार लीजिए,
        सिरे से सिरे तक $\Delta\vect v = \vect v_{i+1} - \vect v_{i-1}$ खींचिए; फिर $\vect a_i \approx \Delta\vect v/(2\,\Delta t)$।
\end{enumerate}
\end{method}

\begin{omfigure}
\begin{tikzpicture}[scale=1.05, font=\small]
  \fill (0,2.05) circle (2pt) (0.8,1.95) circle (2pt)
    (1.6,1.72) circle (2pt) (2.4,1.36) circle (2pt)
    (3.2,0.87) circle (2pt) (4.0,0.25) circle (2pt);
  \node[above] at (0,2.05) {$M_0$};
  \node[above] at (1.6,1.72) {$M_2$};
  \node[above right] at (2.4,1.36) {$M_3$};
  \node[above right] at (3.2,0.87) {$M_4$};
  \draw[-Stealth, omDef, thick] (1.6,1.72) -- ++(1.2,-0.44)
    node[right] {$\vect v_2$};
  \draw[-Stealth, omDef, thick] (3.2,0.87) -- ++(1.2,-0.83)
    node[right] {$\vect v_4$};
  \draw[-Stealth, omDef, dashed] (2.4,1.36) -- ++(1.2,-0.44);
  \draw[-Stealth, omDef] (2.4,1.36) -- ++(1.2,-0.83);
  \draw[-Stealth, omThm, very thick] (3.6,0.92) -- ++(0,-0.39)
    node[right] {$\Delta\vect v$};
  \draw[-Stealth, omThm, very thick] (2.4,1.36) -- ++(0,-0.75)
    node[left] {$\vect a$};
\end{tikzpicture}

{\small फेंकी गई गेंद का \omterm{met:g12:kinematics-2d:chrono}{कालक्रम-चित्र}: जीवाओं से बने $\vect v_2$ और
$\vect
v_4$ को $M_3$ पर उतारा गया है; और उनका अंतर $\Delta\vect v$ $\vect a$ सौंप देता है ---
\omterm{def:g10:universal-gravitation:weight}{ऊर्ध्वाधर}, नीचे की ओर।}
\end{omfigure}

\section{सरल रेखीय गतियाँ}

सरल रेखा पर एक ही अक्ष काफ़ी है: $x(t)$, $v(t) = x'(t)$, $a(t) = v'(t)$, और तीनों चिह्न सहित
संख्याएँ।

\begin{definition}[सरल रेखीय एकसमान गति]\label{def:g12:kinematics-2d:uniform}
कोई गति \emph{सरल रेखीय एकसमान}\index{सरल रेखीय एकसमान गति} तब है जब $\vect v$
अचर सदिश हो: सीधा \omterm{def:g10:relative-motion:trajectory}{गतिपथ}, अचर चाल, $\vect a = \vect 0$। तब $x(t) = v\,t + x_0$।
\end{definition}

\begin{proposition}[सरल रेखा में एकसमान त्वरित गति]\label{prop:g12:kinematics-2d:uarm}
\index{एकसमान त्वरित गति}%
यदि कोई बिंदु $Ox$ के अनुदिश \emph{अचर} त्वरण $a$ के साथ चले और $t = 0$
पर स्थिति $x_0$ से वेग $v_0$ लेकर चले, तो
\[
v(t) = a\,t + v_0,
\qquad
x(t) = \tfrac12\,a\,t^2 + v_0\,t + x_0 .
\]
\end{proposition}

\begin{proof}
$v$ अचर $a$ का प्रतिअवकलज है: $v(t) = a t + C$, और $t = 0$ से $C = v_0$ \omterm{def:g10:inertia:force} जाता
है। उसी तरह $x$ $a t + v_0$ का प्रतिअवकलज है: $x(t) = \tfrac12 a t^2 + v_0 t + C'$, जहाँ $C' = x_0$।
\end{proof}

\begin{omfigure}
\begin{tikzpicture}[font=\small]
\begin{axis}[omaxis, name=ax1, width=8.4cm, height=3.2cm,
    ylabel={$x$ (m)}, xmin=0, xmax=4, ymin=0, ymax=21]
  \addplot[omDef, thick, domain=0:4, samples=40] {0.5*2*x^2 + x};
\end{axis}
\begin{axis}[omaxis, name=ax2, at={(ax1.below south west)},
    anchor=north west, yshift=-3mm, width=8.4cm, height=3.2cm,
    ylabel={$v$ (m/s)}, xmin=0, xmax=4, ymin=0, ymax=10]
  \addplot[omThm, thick, domain=0:4, samples=2] {2*x + 1};
\end{axis}
\begin{axis}[omaxis, at={(ax2.below south west)}, anchor=north west,
    yshift=-3mm, width=8.4cm, height=3.2cm, xlabel={$t$ (s)},
    ylabel={$a$ ($\unit{m/s^2}$)}, xmin=0, xmax=4, ymin=0, ymax=3]
  \addplot[omProp, thick, domain=0:4, samples=2] {2};
\end{axis}
\end{tikzpicture}

{\small \omterm{prop:g12:kinematics-2d:uarm}{एकसमान त्वरित गति} ($a = \qty{2.0}{m/s^2}$, $v_0 = \qty{1.0}{m/s}$, $x_0 = 0$): हर आलेख अपने
ऊपर वाले का अवकलज है --- एक परवलय, उसकी ढाल, और ढाल की ढाल।}
\end{omfigure}

\begin{example}[रुकने की दूरी]\label{ex:g12:kinematics-2d:braking}
$v_0$ चाल वाली कार अचर मंदन $a$ (त्वरण $-a$) से ब्रेक लगाती
है। वह तब रुकती है जब $v(t) = -a t + v_0 = 0$, यानी $t_s = v_0/a$ पर, और तब तक $d = -\tfrac12 a t_s^2 + v_0 t_s =
v_0^2/(2a)$ तय कर चुकी
होती है। सड़क-सुरक्षा की सुर्ख़ी यही वर्ग है: चाल दुगुनी कीजिए और
रुकने की दूरी \emph{चौगुनी} हो जाती है --- $a =
\qty{6.0}{m/s^2}$ के साथ, \qty{50}{km/h} से
\qty{16}{m}, और \qty{100}{km/h} से \qty{64}{m}।
\end{example}

\section{एकसमान वृत्तीय गति और फ्रेने आधार}

\begin{definition}[एकसमान वृत्तीय गति]\label{def:g12:kinematics-2d:ucm}
कोई गति \emph{एकसमान वृत्तीय}\index{एकसमान वृत्तीय गति} तब है जब
\omterm{def:g10:relative-motion:trajectory}{गतिपथ} त्रिज्या $R$ का वृत्त हो और चाल $v$ अचर हो।
\emph{आवर्तकाल}\index{आवर्तकाल!घूर्णन का} $T$ एक चक्कर की अवधि है, और
\emph{आवृत्ति}\index{आवृत्ति!घूर्णन की} $f = 1/T$ (\unit{Hz} में) प्रति \omterm{def:g10:orders-of-magnitude:unit}{सेकंड}
चक्करों की संख्या: $v = 2\pi R/T = 2\pi R f$।
\end{definition}

अचर चाल --- फिर भी वेग मुड़ता जाता है: तो त्वरण \emph{है}। किस ओर, और
कितना बड़ा?

\begin{theorem}[अभिकेंद्र त्वरण]\label{thm:g12:kinematics-2d:centripetal}
\index{अभिकेंद्र त्वरण}%
त्रिज्या $R$ और चाल $v$ वाली \omterm{def:g12:kinematics-2d:ucm}{एकसमान वृत्तीय गति} में
त्वरण --- जिसे \emph{अभिकेंद्र} कहते हैं --- चलते बिंदु से वृत्त के केंद्र की
ओर संकेत करता है, और उसका मानक अचर रहता है:
\[
a = \frac{v^2}{R}.
\]
\end{theorem}

\begin{proof}
अक्षों को वृत्त के केंद्र पर रखिए और $\omega = v/R$ लीजिए, ताकि चाप $v t$ कोण
$\omega t$ बनाए: तब स्थिति $x(t) = R\cos(\omega t)$, $y(t) = R\sin(\omega t)$ है। अवकलन करने पर $\vect v = (-R\omega\sin(\omega t),\, R\omega\cos(\omega t))$, जिसका
मानक अचर $R\omega = v$ है; और दुबारा अवकलन करने पर $\vect a = -\omega^2\,\vect{OM}$: यानी \omterm{def:g12:kinematics-2d:position}{स्थिति
सदिश} के विपरीत --- केंद्र की ओर --- और मानक $\omega^2 R = v^2/R$।
\end{proof}

\begin{omfigure}
\begin{tikzpicture}[scale=1.0, font=\small]
  \draw[thick] (0,0) circle (1.7);
  \fill (0,0) circle (1.5pt); \node[below left] at (0,0) {$O$};
  \draw[dashed, omExo] (0,0) -- (215:1.7)
    node[midway, below right] {$R$};
  \fill (35:1.7) circle (2pt); \node[above right] at (35:1.7) {$M$};
  \draw[-Stealth, omDef, thick] (35:1.7) -- ++(-0.69,0.98)
    node[above] {$\vect v$};
  \draw[-Stealth, omThm, very thick] (35:1.7) -- ++(215:0.95)
    node[below right] {$\vect a$};
\end{tikzpicture}

{\small \omterm{def:g12:kinematics-2d:ucm}{एकसमान वृत्तीय गति}: $\vect v$ स्पर्शी, नियत लंबाई का; $\vect a$
केंद्र की ओर ताका हुआ, नियत लंबाई $v^2/R$ का --- यानी वेग को सदा मोड़ता हुआ,
पर उसे कभी खींचता नहीं।}
\end{omfigure}

\begin{definition}[फ्रेने आधार]\label{def:g12:kinematics-2d:frenet}
वक्र पथ पर चलते बिंदु की वर्तमान स्थिति पर \emph{फ्रेने आधार}\index{फ्रेने
आधार} परस्पर लंब एकांक सदिशों का युग्म $(\vect{u_t}, \vect{u_n})$ है: $\vect{u_t}$ स्पर्शी,
गति के अनुदिश; और $\vect{u_n}$ \omterm{def:g10:refraction:angles}{अभिलंब}, उस त्रिज्या $R$ वाले वृत्त के केंद्र की
ओर जो वक्र से वहाँ सबसे अच्छा बैठता है।
\end{definition}

\begin{proposition}[फ्रेने आधार में त्वरण]\label{prop:g12:kinematics-2d:frenetacc}
स्थानीय त्रिज्या $R$ वाले किसी भी समतल पथ पर, जिसे चाल $v(t)$
से तय किया जा रहा हो,
\[
\vect a = a_t\,\vect{u_t} + a_n\,\vect{u_n},
\qquad a_t = \frac{\dd v}{\dd t},
\qquad a_n = \frac{v^2}{R}:
\]
जहाँ \emph{स्पर्शरेखीय त्वरण}\index{स्पर्शरेखीय त्वरण} $a_t$
चाल बदलता है और \emph{अभिलंब त्वरण}\index{अभिलंब त्वरण} $a_n$ वेग
को मोड़ता है।
\end{proposition}
\admitted

\begin{remark}[उपपत्ति कहाँ रहती है]\label{rem:g12:kinematics-2d:frenetproof}
हमने दोनों सिरे सिद्ध कर लिए हैं: रेखा पर विशुद्ध $a_t$ (\cref{prop:g12:kinematics-2d:uarm}) और वृत्त
पर विशुद्ध $a_n$ (\cref{thm:g12:kinematics-2d:centripetal}); और सामान्य विभाजन वर्ष~1 के खंड में निकाला गया
है।
\end{remark}

\begin{example}[मोड़ में ब्रेक]\label{ex:g12:kinematics-2d:bend}
कोई कार $\dd v/\dd t = \qty{-3.0}{m/s^2}$ से ब्रेक लगाते हुए त्रिज्या $R = \qty{50}{m}$ के मोड़ में $v = \qty{15}{m/s}$ पर
घुसती है: $a_t =
\qty{-3.0}{m/s^2}$, $a_n = 15^2/50 = \qty{4.5}{m/s^2}$, इसलिए $a =
\sqrt{3.0^2 + 4.5^2} \approx \qty{5.4}{m/s^2}$, जो पीछे की ओर \emph{और} मोड़ के
भीतर की ओर संकेत करता है। टायरों को दोनों एक साथ देने पड़ते हैं --- और इसीलिए
दौड़ के चालक मोड़ से \emph{पहले} ब्रेक लगाते हैं।
\end{example}

\section{अभ्यास}

\begin{exercise}[$\star$]\label{exo:g12:kinematics-2d:1}
$x(t) = 3.0\,t$, $y(t) = 4.0\,t$ (\omterm{def:g10:orders-of-magnitude:unit}{मीटर}, \omterm{def:g10:orders-of-magnitude:unit}{सेकंड}) के लिए $\vect v$, चाल, \omterm{def:g10:relative-motion:trajectory}{गतिपथ}
और गति का वर्ग बताइए।
\end{exercise}

\begin{exercise}[$\star$]\label{exo:g12:kinematics-2d:2}
हर \qty{0.10}{s} पर लिए गए \omterm{met:g12:kinematics-2d:chrono}{कालक्रम-चित्र} में बिंदु एक रेखा पर, बराबर
\qty{0.35}{m} की दूरी पर पड़े हैं। गति पहचानिए और चाल निकालिए। खुलते
अंतराल किसका संकेत होते?
\end{exercise}

\begin{exercise}[$\star$]\label{exo:g12:kinematics-2d:3}
$Ox$ के अनुदिश $x(t) = 2.0\,t^2$ के लिए (\omterm{def:g10:orders-of-magnitude:unit}{मीटर}, \omterm{def:g10:orders-of-magnitude:unit}{सेकंड}) $v(t)$ और $a(t)$ निकालिए,
गति का वर्ग बताइए, और $t = \qty{3.0}{s}$ पर चाल निकालिए।
\end{exercise}

\begin{exercise}[$\star$]\label{exo:g12:kinematics-2d:4}
हिंडोले का घोड़ा त्रिज्या \qty{4.5}{m} के वृत्त पर हर \qty{8.0}{s} में एक चक्कर लगाता
है: आवृत्ति, चाल और त्वरण (मानक और दिशा) निकालिए।
\end{exercise}

\begin{exercise}[$\star$]\label{exo:g12:kinematics-2d:5}
कोई $v(t)$ आलेख $t = 0$ पर \qty{2.0}{m/s} से सीधा चढ़कर $t = \qty{4.0}{s}$ पर \qty{10.0}{m/s} हो जाता
है। उससे त्वरण पढ़िए, फिर तय की गई दूरी (आलेख के नीचे का क्षेत्रफल)।
\end{exercise}

\begin{exercise}[$\star\star$]\label{exo:g12:kinematics-2d:6}
कोई तेज़ रफ़्तार रेलगाड़ी त्रिज्या \qty{6.0}{km} का मोड़ \qty{320}{km/h} पर पार
करती है। $a_n$ निकालिए और उसकी तुलना $g$ से कीजिए; तेज़ रफ़्तार
पटरियों को कसे हुए मोड़ों से, चाहे वे पूरी तरह झुके ही क्यों न हों, क्यों
बचना पड़ता है?
\end{exercise}

\begin{exercise}[$\star\star$]\label{exo:g12:kinematics-2d:7}
सीधी पटरी पर चलती ठेली के लिए $x(t) = -1.2\,t^2 + 8.0\,t + 3.0$ है (\omterm{def:g10:orders-of-magnitude:unit}{मीटर}, \omterm{def:g10:orders-of-magnitude:unit}{सेकंड})। $v(t)$ और $a$
निकालिए, फिर पलटने का क्षण और स्थान, और उससे पहले तथा बाद की गति का वर्णन
कीजिए।
\end{exercise}

\begin{exercise}[$\star\star$]\label{exo:g12:kinematics-2d:8}
किसी रेखा के अनुदिश हर $\Delta t = \qty{0.20}{s}$ पर लिया गया \omterm{met:g12:kinematics-2d:chrono}{कालक्रम-चित्र} $x = 0$,
$0.10$, $0.28$, $0.54$, $\qty{0.88}{m}$ देता है। बीच के तीनों बिंदुओं पर वेग का
अनुमान लगाइए (\cref{met:g12:kinematics-2d:chrono}), फिर त्वरण का; और निष्कर्ष लिखिए।
\end{exercise}

\begin{exercise}[$\star\star$]\label{exo:g12:kinematics-2d:9}
चंद्रमा \qty{27.3}{d} में त्रिज्या \qty{3.84e8}{m} के लगभग वृत्त पर पृथ्वी की
परिक्रमा करता है। उसकी चाल और उसका \omterm{thm:g12:kinematics-2d:centripetal}{अभिकेंद्र त्वरण}
निकालिए; और यह जानते हुए कि चंद्रमा लगभग $60$ पृथ्वी-त्रिज्या दूर है
(\cref{ch:g10:universal-gravitation}), दूसरे की तुलना $g$ से कीजिए।
\end{exercise}

\begin{exercise}[$\star\star$]\label{exo:g12:kinematics-2d:10}
त्रिज्या \qty{25}{cm} का कपड़े धोने की मशीन का ढोल $1200$ चक्कर प्रति मिनट से
घूमता है। $f$, $T$, किनारे की चाल और किनारे का त्वरण $g$
के गुणजों में निकालिए। कपड़े चिपके क्यों रह जाते हैं जबकि पानी निकल जाता है?
\end{exercise}

\begin{exercise}[$\star\star$]\label{exo:g12:kinematics-2d:11}
\qty{90}{km/h} पर चलती कार अचर $a = \qty{7.5}{m/s^2}$ से ब्रेक लगाती है: रुकने का समय और ब्रेक-दूरी
निकालिए, और $v(t)$ का रेखाचित्र बनाइए।
\end{exercise}

\begin{exercise}[$\star\star\star$]\label{exo:g12:kinematics-2d:12}
$x(t) = R\cos(\omega t)$, $y(t) = R\sin(\omega t)$ के लिए: $\vect v$ निकालिए; जाँचिए कि चाल अचर $R\omega$
है और $\vect v \perp \vect{OM}$; फिर $\vect a$ निकालिए, दिखाइए कि वह $-\omega^2\,\vect{OM}$ है, और $a = v^2/R$ फिर
से पाइए।
\end{exercise}

\begin{exercise}[$\star\star\star$]\label{exo:g12:kinematics-2d:13}
कोई मोटरसाइकिल \qty{2.5}{m/s^2} से ब्रेक लगाते हुए त्रिज्या \qty{80}{m} के मोड़ में
\qty{20}{m/s} पर घुसती है। $a_t$, $a_n$, $\vect a$ का मानक और गति की दिशा से उसका
कोण निकालिए। टायर कुल मिलाकर \qty{7.0}{m/s^2} तक पकड़ते हैं: क्या सवार सीमा के भीतर
है?
\end{exercise}

\begin{exercise}[$\star\star\star$]\label{exo:g12:kinematics-2d:14}
कोई मेट्रो \omterm{def:g10:relative-motion:frame}{विराम} से \qty{1.2}{m/s^2} पर \qty{20}{m/s} तक त्वरित होती है, \qty{40}{s} तक उसी चाल
पर चलती है, फिर \qty{1.5}{m/s^2} से ब्रेक लगाकर रुक जाती है। हर चरण की अवधि और लंबाई,
कुल दूरी और \omterm{def:g10:relative-motion:speed}{औसत चाल} निकालिए। $v(t)$ का रेखाचित्र बनाइए और दूरी को
क्षेत्रफल के रूप में जाँचिए।
\end{exercise}

\begin{exercise}[$\star\star\star$]\label{exo:g12:kinematics-2d:15}
कोई भूस्थिर \omterm{def:g10:universal-gravitation:satellite}{उपग्रह} पृथ्वी के केंद्र से त्रिज्या \qty{4.22e7}{m} पर
\omterm{def:g12:kinematics-2d:ucm}{आवर्तकाल} \qty{86164}{s} के साथ चक्कर लगाता है। उसकी चाल और उसका
\omterm{thm:g12:kinematics-2d:centripetal}{अभिकेंद्र त्वरण} निकालिए। उसे देता कौन है, और \omterm{def:g10:universal-gravitation:satellite}{उपग्रह} भूमध्य
रेखा के एक ही नियत बिंदु के ऊपर क्यों टँगा रहता है (\cref{ch:g12:satellites-kepler})?
\end{exercise}

\section{समस्या: दुर्घटना की रिपोर्ट}

\begin{problem}[{सप्ताहांत समस्या --- दुर्घटना की रिपोर्ट: फ़्रेम-दर-फ़्रेम
पढ़ी गई डैशकैम की फ़िल्म, घिसटन के निशानों से फिर खड़ा किया गया ब्रेक का नियम,
पकड़ के बारे में गवाही देता एक गोल चक्कर, और किलोमीटर प्रति घंटे में सुनाया
गया फ़ैसला}]
\label{pb:g12:kinematics-2d:1}
गली से निकलती एक वैन से कार टकरा जाती है; चोट किसी को नहीं आई, पर बीमा कंपनियाँ
आपस में सहमत नहीं। जाँचकर्ता के पास ऊपर से लिया गया कैमरा-फ़ुटेज (हर \qty{0.50}{s}
पर एक फ़्रेम), नापे गए घिसटन के निशान और कार की डैशकैम है। चाल-सीमा
\qty{50}{km/h} है।

\medskip
\noindent\textbf{भाग I --- फ़िल्म पढ़ना।}
किसी भी प्रतिक्रिया से पहले के फ़्रेम, सीधी सड़क के अनुदिश, ये देते हैं:
\begin{center}\small
\begin{tabular}{lccccc}
$t$ (\unit{s}) & 0.00 & 0.50 & 1.00 & 1.50 & 2.00\\
$x$ (\unit{m}) & 0.0 & 9.0 & 18.0 & 27.0 & 36.0
\end{tabular}
\end{center}
\begin{enumerate}
  \item रिपोर्ट को सबसे पहले तंत्र और अक्ष क्यों बताने चाहिए? यहाँ कौन-से
        लिए गए हैं?
  \item चारों अंतरालों में से हर एक पर औसत वेग निकालिए।
  \item यह गति किस वर्ग की है? सारणी से कारण दीजिए।
  \item कार की चाल $v_0$ \unit{m/s} और \unit{km/h} में दीजिए।
  \item इन दो सेकंडों में त्वरण क्या है?
\end{enumerate}

\medskip
\noindent\textbf{भाग II --- घिसटन के निशान।}
इस सूखी डामर पर टायर-परीक्षण अचर ब्रेक-मंदन $a = \qty{6.0}{m/s^2}$ देते हैं; और घिसटन के
निशान $d_b = \qty{27}{m}$ लंबे हैं। ब्रेक लगने के आरंभ पर $t = 0$, $x = 0$ लीजिए, और
अनजान चाल $v_B$।
\begin{enumerate}[resume]
  \item प्रतिअवकलजों से ब्रेक लगने के दौरान $v(t)$ और $x(t)$ लिखिए।
  \item रुकने का समय $t_s$ लिखिए; दिखाइए कि $d_b = v_B^2/(2a)$।
  \item घिसटन के निशानों से $v_B$ निकालिए। क्या यह भाग~I से संगत है?
  \item प्रतिक्रिया-समय \qty{1.0}{s} है: प्रतिक्रिया-दूरी और रुकने की कुल दूरी
        निकालिए।
  \item प्रश्न 7 से समझाइए कि चाल दुगुनी करने पर ब्रेक-दूरी चौगुनी
        क्यों हो जाती है; और ठीक \qty{50}{km/h} पर $d_b$ निकालिए।
\end{enumerate}

\medskip
\noindent\textbf{भाग III --- गोल चक्कर की गवाही।}
इससे पहले डैशकैम कार को त्रिज्या $R = \qty{12.5}{m}$ वाली गोल चक्कर की गली में अचर
चाल से घूमते दिखाती है, और चौथाई चक्कर \qty{2.5}{s} में। सड़क अधिक से अधिक
$a_n = \mu g$ दे सकती है, जहाँ $\mu = 0.70$।
\begin{enumerate}[resume]
  \item वहाँ चाल अचर है --- क्या वेग भी? क्या कार त्वरित हो रही है?
        समझाइए।
  \item गोल चक्कर पर कार की चाल निकालिए।
  \item उसका \omterm{thm:g12:kinematics-2d:centripetal}{अभिकेंद्र त्वरण} (मानक और दिशा) निकालिए।
  \item बिना फिसले अधिकतम चाल निकालिए। क्या कार उसके भीतर थी?
  \item उस चाल पर पूरे चक्कर का \omterm{def:g12:kinematics-2d:ucm}{आवर्तकाल} और
        आवृत्ति दीजिए।
\end{enumerate}

\medskip
\noindent\textbf{भाग IV --- फ़ैसला।}
जिस क्षण चालक ने प्रतिक्रिया शुरू की, उस समय वैन कार से $D = \qty{41}{m}$ आगे प्रकट हुई
थी।
\begin{enumerate}[resume]
  \item ब्रेक लगने से पहले की चाल सीमा से कितनी अधिक थी, \unit{km/h} में
        और प्रतिशत में?
  \item टक्कर के क्षण कार की चाल निकालिए।
  \item ठीक \qty{50}{km/h} पर (वही प्रतिक्रिया, वही ब्रेक) रुकने की गणना दोहराइए:
        क्या कार रुक जाती, और कितने हाशिये से?
  \item रुकने की दूरी का सामान्य सूत्र $d(v)$ दीजिए; उसमें चाल दो
        बार क्यों आती है --- एक बार रैखिक रूप में, एक बार वर्ग में?
  \item फ़ैसला, संख्याओं समेत एक वाक्य में: चालक किस चाल पर था, और
        क्या सीमा पर चलते हुए यह दुर्घटना होती?
\end{enumerate}
\end{problem}
